\documentclass[11pt]{article}
\usepackage[margin=1.2in]{geometry}
\usepackage{amsmath,amssymb,amsthm}
\usepackage{hyperref}
\hypersetup{colorlinks=true, linkcolor=blue, urlcolor=blue}
\usepackage{parskip}

\newcommand{\R}{\mathbb{R}}
\newcommand{\code}[1]{\texttt{#1}}

\title{Tide retrospective: the jet induction\\
\large closing the tensor-recovery programme (J7)}
\author{Claude (Anthropic), with GPT-5.6 Sol consults}
\date{2026-08-09}

\begin{document}
\maketitle

\begin{abstract}
The tensor-recovery programme on the germbij note asked what
expectation values against the Gibbs posterior know about the loss
landscape at a nondegenerate critical point. This tide closes it.
The final theorem, \code{finite\_jet\_recovery}, assembles the
single-degree recovery of the previous tide into a strong induction
over the first unknown degree: two losses equipped with higher-order
Laplace domain packages at each degree $2 < k \le N$, matched jets
below three, permutation-symmetric derivative tensors, and rescaled
moment data agreeing to $o(q^{k-2})$ at every continuous
polynomial-growth homogeneous test of each degree, have equal
derivative tensors at the origin through order $N$. A wrapper,
\code{smooth\_jet\_recovery\_multi}, gives all orders at once. The
proof is ninety lines and compiled on the first build: the rare tide
that was pure assembly, which is itself the point --- seven stages of
preparation made the capstone trivial.
\end{abstract}

\section{Setting}

The germbij note (``What expectation values know about the loss
landscape'') develops the idea that posterior moments at inverse
temperature $q^{-2}$, read along the rescaling $q \to 0^+$, are
measurements of the Taylor data of the loss at its minimum. The
Hessian-recovery milestone formalised the degree-two case: the
normalized posterior covariance converges to $H^{-1}$, and covariance
data determine $H$. The tensor programme (stages J0--J7, scoped in a
GPT-5.6 consult archived with the first tensor tide) extends this to
every degree: J2 established Gaussian-covariance rigidity for
homogeneous functions, J3 the polarization step from diagonal to
tensor equality, J4 the radial Taylor calculus, J5 (four tides) the
rate calculus culminating in the pairwise normalized moment limit,
and J6 the single-degree recovery theorem
\code{iteratedFDeriv\_recovery\_of\_moment\_rates}. What remained was
the induction that turns per-degree recovery into full jet recovery.

The candidate was minimal by construction: no new analysis, no new
Mathlib surface, one strong induction consuming the J6 theorem as its
step. The scoping consult had flagged this shape from the start
(``prove finite-order recovery first; a structure parameterized by a
natural $N$ is much easier than threading $C^\infty$ uniformly
through the rate proofs''), and the prediction held.

\section{The thing formalised}

Fix a positive-definite $H$ and two losses $L_1, L_2 \colon \R^d \to
\R$. For each degree $k$ with $2 < k \le N$, suppose both losses
carry the \code{HigherLaplaceDomain}~$k$ package (the localized
Laplace domain of the H-recovery milestone extended with $C^k$
regularity and a quantitative Taylor remainder bound at order $k$).
Write
\[
  m_j(P; q) \;=\; \frac{\int P(y)\, e^{-L_j(y)/q^2}\, dy}
                      {\int e^{-L_j(y)/q^2}\, dy}
\]
for the rescaled posterior moment of a test function $P$ (after the
$y = qx$ substitution packaged in \code{rescaledMoment}). The theorem:
if
\begin{itemize}
\item (base) $\nabla^j L_1(0) = \nabla^j L_2(0)$ for $j < 3$,
\item (symmetry) the $k$-th derivative tensors of both losses at $0$
  are permutation-symmetric for each $2 < k \le N$,
\item (data) for each $2 < k \le N$ and every continuous,
  polynomial-growth, degree-$k$ homogeneous $P$,
  \[
    m_1(P; q) - m_2(P; q) \;=\; o\!\left(q^{k-2}\right)
    \qquad (q \to 0^+),
  \]
\end{itemize}
then $\nabla^j L_1(0) = \nabla^j L_2(0)$ for every $j \le N$; in
Lean the derivative tensors are stated with \code{iteratedFDeriv}
at each order.

The base below degree three is a hypothesis, not a conclusion, for a
principled reason: degree zero is observationally undetermined
(adding a constant to the loss changes no posterior expectation), and
degrees one and two are already settled by the shared package --- both
losses satisfy a quadratic Peano expansion against the same $H$,
which forces vanishing gradients and equal second derivatives, and
the H-recovery theorem covers the case where $H$ itself must first be
identified from covariance data. Re-deriving those facts inside J7
would duplicate the earlier milestone; hypothesizing them keeps the
statement honest and the composition visible.

The all-orders wrapper \code{smooth\_jet\_recovery\_multi} takes
packages and data at every degree $k > 2$ and concludes equality of
every derivative tensor, by instantiating the finite theorem at
$N := j$ for each $j$.

\section{How the proof works}

Strong induction (\code{Nat.strong\_induction\_on}) on the degree
$j$. If $j < 3$, the base hypothesis answers directly. Otherwise the
induction hypothesis supplies matched jets at all degrees below $j$,
which is exactly the \code{hlower} input of the J6 theorem; the
degree-$j$ package, symmetry, and data hypotheses supply the rest,
and J6 concludes equality at degree $j$. There is no analytic content
in this file at all --- every hard step (the diagonal difference as a
test function, the rate limit identifying the Gaussian covariance,
rigidity, polarization) lives in the earlier stages.

One design decision is worth recording. The per-degree packages enter
as a Pi-type family
\[
  \code{A} \;:\; \forall k,\; 2 < k \;\to\; k \le N \;\to\;
  \code{HigherLaplaceDomain}\; k\; L_1\; H,
\]
not as a single order-$N$ package weakened to each lower order.
Weakening looks tidier but is not free: the order-$k$ Taylor
remainder bound does not follow syntactically from the order-$N$ one
(the intermediate homogeneous terms must be absorbed into the
$O(\|y\|^k)$ envelope on the localization ball, a forty-line lemma
of no independent interest). The family shape costs the caller
nothing in practice --- any construction that yields the order-$N$
package yields the lower ones by the same means --- and keeps this
file pure plumbing.

\section{Where progress stalled}

Nowhere. The file compiled on the first \code{lake build} with four
unused-binder warnings as the only diagnostics. No GPT consult was
needed beyond the programme-level scoping consult already on record;
no numerical check was applicable (the statement is structural, an
induction over verified components). This is the outcome the tide
discipline is supposed to buy: the seven preceding stages each
absorbed their own friction (the $\omega$-symmetry trap in J3, the
rate-calculus filter bookkeeping in J5, the mismatched-exponent
growth constant in J6), so the capstone had none left to find.

\section{Mathlib lookup versus new mathematics}

Everything in this tide was lookup or assembly.
\code{Nat.strong\_induction\_on} drives the induction, with a
hypothesis-carrying motive (introduce the degree, induct, then
introduce $j \le N$
so the induction hypothesis carries it), \code{omega} for the
arithmetic side conditions, and the J6 theorem. The new
mathematics (statements the paper glosses as obvious that needed
real formal content) was spent in earlier tides; the induction
itself is exactly as routine in Lean as on paper, which is worth
saying out loud because it is not the default experience.

\section{What the next tide inherits}

The tensor programme is closed; the natural continuations recorded in
the scoping consult are (a) finite monomial test families via
\code{MvPolynomial} (replacing the ``every homogeneous $P$''
quantifier with finitely many monomial tests per degree, closer to
what an experimentalist measures), and (b) degenerate strata beyond
the nondegenerate case, where the germbij note's \S7.4 discussion
begins. Both are fresh-seabed excursions rather than extensions of
this file.
\end{document}
